Rozdział ten opisuje architekturę środowiska badawczego oraz szczegółową metodykę przeprowadzenia testów porównawczych wytypowanych baz danych typu klucz-wartość. Przedstawiono w nim charakterystykę sprzętowo-programową stanowiska pomiarowego, konfigurację środowiska kontenerowego, a także założenia i scenariusze testowe zmierzające do weryfikacji wydajności badanych silników.

\section{Wykorzystane technologie i narzędzia}
W celu zapewnienia powtarzalności badań i eliminacji wpływu środowiska uruchomieniowego, badania oparto na technologiach konteneryzacji oraz automatyzacji testów:
\begin{itemize}
    \item \textbf{Docker i Docker Compose:} wykorzystane do stworzenia odizolowanych kontenerów dla każdego z badanych systemów bazodanowych. Pozwala to na precyzyjne przydzielanie zasobów sprzętowych, takich jak procesor czy pamięć RAM, i łatwe resetowanie stanu baz przed każdym testem.
    \item \textbf{Python (wersja 3.14):} główny język aplikacji benchmarkującej (Benchmark Runner). Wybór ten wynika z dostępności bibliotek klienckich dla wszystkich badanych baz danych oraz łatwości tworzenia skryptów sterujących testami.
    \item \textbf{Systemy operacyjne:} środowisko wirtualizacji WSL 2 (Windows Subsystem for Linux 2) \cite{wslperf2021} pod kontrolą systemu Windows oraz natywny system Linux. W obu wariantach jako bazowy system operacyjny wykorzystano dystrybucję Ubuntu 24.04 LTS, co pozwala na zachowanie tożsamego środowiska pakietów w obu środowiskach testowych.
\end{itemize}

\section{Konfiguracja sprzętowa}
Wszystkie eksperymenty badawcze przeprowadzono na jednolitej platformie sprzętowej o następującej konfiguracji:
\begin{itemize}
    \item \textbf{Procesor:} ośmiordzeniowy i szesnastowątkowy procesor AMD Ryzen 7 5700U (8 rdzeni fizycznych z obsługą SMT), o taktowaniu bazowym 1,8 GHz (do 4,3 GHz w trybie Boost) oraz 8 MB pamięci podręcznej L3.
    \item \textbf{Pamięć RAM:} 16 GB pamięci LPDDR4X działającej w trybie dwukanałowym.
    \item \textbf{Dysk twardy:} nośnik SSD M.2 PCIe NVMe o pojemności 1 TB.
\end{itemize}
W celu zapewnienia powtarzalności pomiarów w środowisku WSL 2, narzucono sztywne limity zasobów obliczeniowych i pamięciowych przydzielonych dla maszyny wirtualnej. Konfigurację tę zdefiniowano w pliku \texttt{.wslconfig} w katalogu profilu użytkownika systemu operacyjnego Windows, przydzielając maszynie wirtualnej limit 8~GB pamięci RAM (odpowiadający połowie fizycznej pamięci operacyjnej platformy sprzętowej) oraz dostęp do wszystkich 16~wątków logicznych procesora. Szczegółową konfigurację tego pliku przedstawiono na listingu~\ref{lst:wslconfig}.

\begin{lstlisting}[label={lst:wslconfig},caption={Konfiguracja zasobów środowiska WSL 2 w pliku .wslconfig}]
[wsl2]
memory=8GB
processors=16
\end{lstlisting}

\section{Konfiguracja sieci kontenerów}
Do komunikacji pomiędzy aplikacją testową napisaną w języku Python, a kontenerami bazodanowymi wykorzystano dedykowaną, izolowaną sieć wirtualną utworzoną w silniku Docker (bridge network). Zapewniło to pełną izolację komunikacji benchmarkowej od zewnętrznego ruchu sieciowego hosta, gwarantując stabilne i powtarzalne parametry transmisji sieciowej podczas wszystkich serii pomiarowych.

\section{Struktura repozytorium i organizacja kodu testowego}
Kod testowy został podzielony na następujące moduły:
\begin{itemize}
    \item \textbf{Skrypt uruchomieniowy (Benchmark Launcher):} skrypt powłoki systemowej odpowiedzialny za równoległe uruchamianie i koordynację zadanej liczby niezależnych procesów klienckich bezpośrednio z poziomu terminala.
    \item \textbf{Jednostkowy proces kliencki (Client Runner):} samodzielna instancja programu w języku Python realizująca operacje na bazie danych zgodnie z określonym profilem obciążenia oraz zapisująca własne wyniki cząstkowe w formacie JSON.
    \item \textbf{Adaptery bazodanowe (Database Wrappers):} skrypty implementujące jednolite operacje odczytu, zapisu, aktualizacji i usuwania dla każdej z baz danych.
    \item \textbf{Moduł agregujący:} skrypt odpowiedzialny za konsolidację wyników cząstkowych zebranych przez poszczególne procesy klienckie oraz generowanie wykresów zbiorczych przy użyciu biblioteki \texttt{matplotlib}.
\end{itemize}

\section{Założenia testowe i limity zasobów}
W celu zbadania skalowalności pionowej badanych baz danych, każdy z kontenerów bazodanowych był uruchamiany z precyzyjnie określonymi limitami zasobów procesora. Zastosowanie technologii Docker umożliwiło precyzyjną izolację zasobową z wykorzystaniem mechanizmu grup kontrolnych systemu Linux (cgroups). Zdefiniowano trzy główne konfiguracje limitów CPU dla serwerów bazodanowych:
\begin{itemize}
    \item Limit wynoszący 1 rdzeń CPU (pełna izolacja jednowątkowa serwera).
    \item Limit wynoszący 2 rdzenie CPU.
    \item Limit wynoszący 4 rdzenie CPU (zaangażowanie współbieżności serwera).
\end{itemize}
W ostatecznej konfiguracji badawczej nie stosowano sztucznych ograniczeń pamięci RAM dla kontenerów bazodanowych. Pozwoliło to silnikom bazodanowym na pełne wykorzystanie pamięci operacyjnej przydzielonej do środowiska (WSL 2/Linux) na potrzeby buforowania danych, indeksów i struktur pamięci podręcznej. Dzięki temu eliminowane jest zjawisko sztucznego wypychania danych na dysk (ang. \emph{eviction}) oraz częstych błędów strony, co umożliwia bezpośrednie badanie czystej wydajności algorytmów i mechanizmów synchronizacji silników przy zróżnicowanej liczbie rdzeni CPU.

Ważnym aspektem metodologicznym przy testowaniu wydajności na jednej fizycznej maszynie jest zjawisko głodzenia procesora (ang. \emph{CPU starvation}) oraz rywalizacja o zasoby obliczeniowe. Ponieważ zarówno aplikacja benchmarkująca (klient), jak i badany silnik bazy danych (serwer) działają na tej samej platformie sprzętowej, jednoczesne pełne obciążenie wszystkich wątków przez klienta mogłoby uniemożliwić prawidłowe działanie serwera. W takim scenariuszu zmierzona wydajność nie odzwierciedlałaby wewnętrznych cech baz danych (takich jak rywalizacja o blokady czy narzuty sieciowe), lecz sprawność szeregowania zadań przez planistę systemowego, w szczególności domyślny dla jądra Linux Completely Fair Scheduler (CFS) \cite{love2010}. W celu wyeliminowania tego problemu, nałożenie limitów na liczbę rdzeni przydzielonych dla kontenera bazy (od 1 do 4 rdzeni) gwarantuje, że kontener serwera bazodanowego nie zajmuje całości dostępnych zasobów obliczeniowych, pozostawiając wystarczającą część rdzeni dla procesów klienckich aplikacji benchmarkującej.

\section{Zestaw operacji testowych}

Badania obejmują zarówno podstawowe (CRUD), jak i zaawansowane operacje na danych, co pozwala na kompleksową ocenę sprawności silników w różnych scenariuszach użytkowych:
\begin{itemize}
    \item \textbf{Insert (zapis prosty):} wstawianie par klucz-wartość do bazy danych. Operacja ta pozwala na ocenę wydajności zapisu sekwencyjnego bądź indeksowania nowych kluczy.
    \item \textbf{Read (odczyt punktowy):} pobieranie wartości na podstawie losowo wybranego klucza z zakresu wcześniej wprowadzonych danych.
    \item \textbf{Update (modyfikacja):} aktualizacja wartości (nadpisanie rekordu) pod istniejącym już kluczem.
    \item \textbf{Delete (usuwanie):} usuwanie rekordów z bazy danych na podstawie klucza.
    \item \textbf{Mieszane obciążenie (Mixed Workload):} scenariusz łączący operacje odczytu (GET) oraz zapisu (SET) na tym samym kluczu w ramach jednej pętli testowej (z następowym usunięciem klucza). Testy te realizowane są w trzech proporcjach:
    \begin{itemize}
        \item 90\% odczytów i 10\% zapisów (profil zdominowany przez odczyt).
        \item 50\% odczytów i 50\% zapisów (profil zbalansowany).
        \item 10\% odczytów i 90\% zapisów (profil zdominowany przez zapis).
    \end{itemize}
    Ze względu na złożoność i narzut czasowy testów mieszanych, wolumen operacji został ograniczony do stałej liczby 10~000 iteracji na proces kliencki. Większa liczba powtórzeń nie wnosi dodatkowej wartości analitycznej, a jedynie wydłuża czas trwania testu.
    \item \textbf{Kolejka (Queue):} test symulujący działanie kolejki komunikatów typu FIFO (ang. \emph{First-In, First-Out}). Operacja ta składa się z dwóch faz:
    \begin{itemize}
        \item W pierwszej fazie następuje przeplatana operacja dodawania (PUSH) oraz pobierania (POP) elementów w stosunku 2:1 (dwa dodania na jedno pobranie), co prowadzi do kontrolowanego wzrostu rozmiaru kolejki.
        \item W drugiej fazie pobierane są pozostałe w kolejce elementy aż do jej całkowitego opróżnienia.
    \end{itemize}
    Dla baz serwerowych o bogatej strukturze typów danych (np. Redis) operacja ta jest realizowana natywnie przy użyciu wbudowanych instrukcji listowych (np. \texttt{rpush} i \texttt{lpop}). Dla pozostałych systemów operacja ta jest emulowana programowo po stronie aplikacji w języku Python.
    \item \textbf{Zaawansowane operacje na dokumentach JSON:} testy te badają wydajność operacji na strukturach półstrukturalnych (obiektach JSON zawierających klucze typu \texttt{payload}, \texttt{counter} oraz tablicę \texttt{tags}). Obejmują one następujący cykl:
    \begin{itemize}
        \item \textbf{JSON Insert:} zapisanie pełnego dokumentu JSON.
        \item \textbf{JSON Read:} odczyt całego zapisanego dokumentu.
        \item \textbf{JSON Read Partial:} odczyt jedynie określonego wycinka dokumentu (pola \texttt{payload}).
        \item \textbf{JSON Update Partial:} aktualizacja wyłącznie pola \texttt{payload} wewnątrz dokumentu.
        \item \textbf{JSON Increment:} inkrementacja pola numerycznego \texttt{counter} wewnątrz dokumentu.
        \item \textbf{JSON Delete:} usunięcie całego dokumentu z bazy.
    \end{itemize}
    Podobnie jak w przypadku kolejki, systemy wspierające natywnie dokumenty JSON (np. Redis z modułem RedisJSON, Couchbase, Tarantool) wykonują te modyfikacje po stronie serwera. Systemy prostych kluczy-wartości (np. Memcached, LevelDB, RocksDB, BerkeleyDB, FoundationDB) realizują te operacje poprzez emulację kliencką (pobranie surowej wartości tekstowej, deserializacja do obiektu słownika w języku Python, modyfikacja odpowiedniego pola, ponowna serializacja do formatu JSON i nadpisanie całego rekordu w bazie).
\end{itemize}

\section{Zestaw danych testowych i ograniczenia}
Do testów wykorzystano zestawy danych różniące się rozmiarem pojedynczego rekordu oraz ich łączną liczbą:
\begin{itemize}
    \item \textbf{Rozmiar rekordu (wartości):}
    \begin{itemize}
        \item Mały: 128 bajtów (odpowiednik prostych parametrów i identyfikatorów sesji).
        \item Średni: 4 KB.
    \end{itemize}
    \item \textbf{Wolumen danych (liczba operacji):} testy CRUD przeprowadzono dla trzech wolumenów: 10~000 operacji (szybka weryfikacja), 100~000 operacji oraz 1~000~000 operacji (pełne wysycenie i stabilny pomiar wydajności).
\end{itemize}

W metodyce badań wprowadzono istotne ograniczenie systemowe dotyczące testowania scenariuszy zaawansowanych (kolejki oraz operacji JSON) dla średniego rozmiaru rekordu (4 KB). Testy te są uruchamiane \textbf{wyłącznie} dla małego rozmiaru rekordu (128 bajtów). Ograniczenie to wynika z faktu, że dla systemów emulujących te struktury po stronie klienta (np. Memcached, LevelDB, RocksDB), narzut związany z ciągłym przesyłaniem przez sieć, parsowaniem oraz serializacją obiektów o rozmiarze 4 KB w jednowątkowym interpreterze języka Python całkowicie dominuje czas wykonania zapytania. W efekcie zmierzona wydajność odzwierciedlałaby wydajność biblioteki serializującej języka Python, a nie rzeczywistą przepustowość silnika bazy danych. Dodatkowo niektóre systemy (np. Memcached) posiadają wbudowane limity i ograniczenia rozmiaru przesyłanych kluczy lub specyficzne wymagania protokołu utrudniające stabilne przesyłanie większych struktur emulowanych. Szczególnym przypadkiem jest system FoundationDB, który narzuca sztywne i restrykcyjne limity rozmiarów: maksymalny rozmiar klucza wynosi 10~KB, a wartości 100~KB. W scenariuszu emulacji kolejki FIFO (gdzie stan kolejki jest przechowywany pod pojedynczym kluczem), limit 100~KB dla wartości stanowiłby krytyczną barierę uniemożliwiającą obsługę i skalowanie kolejki przy większej liczbie komunikatów. Z tego względu w adapterze testowym FoundationDB zrezygnowano z przeprowadzenia testu kolejki (kategoria \texttt{queue} pozostała pusta).

Dla scenariuszy CRUD ze średnim rozmiarem rekordu (4 KB) pominięto wolumen 1~000~000 operacji. Wynika to z faktu, iż przy 8 równoległych procesach klienckich łączna liczba unikalnych kluczy przechowywanych w bazie wynosi wówczas $8 \times 1\,000\,000 = 8\,000\,000$ rekordów o rozmiarze 4 KB, co odpowiada zbiorowi danych o łącznej objętości \textbf{32 GB}. Obciążenie to przekracza dostępną pamięć operacyjną platformy testowej i powoduje nadmierne ciśnienie pamięciowe (ang.~\emph{memory pressure}) oraz intensywną wymianę danych z dyskiem, co zafałszowałoby wyniki pomiaru wydajności silnika bazy danych.

Istotnym aspektem metodologicznym jest sposób organizacji przestrzeni kluczy w scenariuszu wieloprocesowym. Każdy z niezależnych procesów klienckich operuje na \textbf{wyłącznie własnym, rozłącznym zbiorze kluczy} (zakresy kluczy są rozdzielone pomiędzy procesy bez nakładania się). Oznacza to, że łączna liczba rekordów przechowywanych jednocześnie w bazie danych rośnie \textbf{proporcjonalnie} do liczby procesów klienckich. Tym samym zużycie pamięci RAM oraz przestrzeni dyskowej przez serwer bazy danych jest wprost zależne od liczby procesów klienckich, przy czym rzeczywista wielkość tego wzrostu zależy od rodzaju wykonywanych operacji (operacje zapisu generują nowe rekordy, odczytu nie powiększają zbioru danych, a usuwania go redukują).


\subsection{Zakres testów i specyficzne ograniczenia silników}
W celu usystematyzowania różnic w zakresie przeprowadzonych testów oraz fizycznych ograniczeń poszczególnych systemów bazodanowych, w tabeli \ref{tab:zakres_i_ograniczenia} zestawiono stopień obsługi współbieżności, charakterystykę realizacji scenariuszy zaawansowanych oraz kluczowe ograniczenia rozmiarowe i architektoniczne nałożone przez poszczególne silniki.

\begin{table}[H]
\centering
\resizebox{\textwidth}{!}{%
\begin{tabular}{|l|>{\raggedright\arraybackslash}p{3.5cm}|>{\raggedright\arraybackslash}p{4.0cm}|>{\raggedright\arraybackslash}p{4.0cm}|>{\raggedright\arraybackslash}p{5.5cm}|}
\hline
\textbf{Baza danych} & \textbf{Dostęp wieloprocesowy} & \textbf{Scenariusz JSON} & \textbf{Scenariusz kolejki FIFO} & \textbf{Główne ograniczenia silnika i rozmiaru} \\ \hline
Redis & Pełny (1-8 procesów) & Natywne wsparcie (zapis / odczyt / modyfikacje) & Natywne wsparcie (struktura LIST) & Brak krytycznych ograniczeń w ramach zakresu testowego. \\ \hline
Memcached & Pełny (1-8 procesów) & Emulacja kliencka (wyłącznie dla rekordu 128 B) & Emulacja kliencka (wyłącznie dla rekordu 128 B) & Domyślny limit rekordu 1~MB zwiększony programowo do 2~MB. \\ \hline
LevelDB & Wyłącznie 1 proces (blokada LOCK) & Emulacja kliencka (wyłącznie dla rekordu 128 B) & Emulacja kliencka (wyłącznie dla rekordu 128 B) & Brak dostępu sieciowego, wyłączna blokada katalogu bazy danych. \\ \hline
RocksDB & Wyłącznie 1 proces (blokada LOCK) & Emulacja kliencka (wyłącznie dla rekordu 128 B) & Emulacja kliencka (wyłącznie dla rekordu 128 B) & Brak dostępu sieciowego, wyłączna blokada katalogu bazy danych. \\ \hline
Couchbase & Pełny (1-8 procesów) & Natywne wsparcie (struktura dokumentowa) & Emulacja kliencka (wyłącznie dla rekordu 128 B) & Maksymalny rozmiar pojedynczego dokumentu (klucza/wartości) wynosi 20~MB. \\ \hline
Tarantool & Pełny (1-8 procesów) & Natywne wsparcie (Lua/przestrzenie \texttt{map}) & Natywne wsparcie (Lua/przestrzeń indeksowana) & Brak krytycznych ograniczeń w ramach zakresu testowego. \\ \hline
FoundationDB & Pełny (1-8 procesów) & Emulacja kliencka (wyłącznie dla rekordu 128 B) & Pominięty (brak możliwości testowania) & Sztywne limity rozmiaru: maksymalna wartość 100~KB. \\ \hline
BerkeleyDB & Pełny (1-8 procesów, DBEnv) & Emulacja kliencka (wyłącznie dla rekordu 128 B) & Emulacja kliencka (wyłącznie dla rekordu 128 B) & Narzut na synchronizację międzyprocesową (blokady środowiska). \\ \hline
\end{tabular}%
}
\caption{Zestawienie zakresu testowego oraz ograniczeń badanych systemów bazodanowych}
\label{tab:zakres_i_ograniczenia}
\end{table}

\section{Scenariusze eksperymentów i współbieżność}
Eksperymenty podzielono na dwa główne scenariusze różniące się stopniem równoległości zapytań:
\begin{itemize}
    \item \textbf{Scenariusz jednowątkowy (Single-Client):} określenie bazowej wydajności każdego z systemów bez narzutu na współbieżność i synchronizację.
    \item \textbf{Scenariusz wieloprocesowy (Multi-Client):} uruchomienie benchmarku na 1, 2, 4 oraz 8 równoległych procesach klienckich.
\end{itemize}
Z uwagi na obecność mechanizmu GIL (ang. \emph{Global Interpreter Lock}) w standardowym interpreterze języka Python (CPython), zastosowanie klasycznych wątków (biblioteka \texttt{threading}) uniemożliwiłoby rzeczywiste równoległe wykonywanie kodu Pythona na wielu rdzeniach procesora. Zdecydowano się również na rezygnację z wbudowanej biblioteki \texttt{multiprocessing} w ramach jednego programu nadrzędnego. Wynika to z faktu, iż nie wszystkie biblioteki klienckie oraz adaptery dedykowane dla badanych baz danych w pełni i stabilnie wspierają współdzielenie zasobów w modelu wieloprocesowości Pythona (często dochodzi do konfliktów przy kopiowaniu gniazd sieciowych bądź deskryptorów plików w procesach potomnych). Aby wyeliminować te ograniczenia i zapewnić pełną, bezpieczną współbieżność po stronie klienta, benchmark zaimplementowano poprzez uruchamianie wielu niezależnych procesów Pythona bezpośrednio z poziomu terminala systemowego. Każdy proces kliencki jest uruchamiany jako osobne, autonomiczne zadanie systemowe posiadające własną przestrzeń adresową oraz niezależny interpreter, co pozwala na generowanie maksymalnego obciążenia serwera bazy danych bez ryzyka błędów sterowników klienckich.

\section{Wizualizacja wyników}
W tym scenariuszu kluczowym celem jest analiza skalowania wydajności systemów wraz ze wzrostem liczby klientów, przy czym maksymalny stopień współbieżności ograniczono do 8 procesów klienckich. Wartość ta odpowiada liczbie fizycznych rdzeni procesora AMD Ryzen 7 5700U (posiadającego 8 rdzeni fizycznych z obsługą 16 wątków logicznych w technologii SMT). Ograniczenie to sprawia, że łączna liczba aktywnych procesów klienckich nie przekracza liczby dostępnych rdzeni fizycznych platformy sprzętowej. Pozwala to planiście zadań systemu operacyjnego na przydzielenie każdemu procesowi klienckiemu osobnego rdzenia fizycznego. Zapobiega to konieczności szeregowania wielu konkurencyjnych procesów na tym samym rdzeniu fizycznym (np. w ramach technologii SMT), co prowadziłoby do rywalizacji o jednostki wykonawcze (ALU, FPU) oraz pamięci podręczne L1/L2 \cite{love2010, kleppmann2017}. Dzięki temu zmierzona wydajność i czasy opóźnień lepiej odzwierciedlają rzeczywistą skalowalność samych silników baz danych, minimalizując narzuty planisty systemowego, częste przełączanie kontekstu (context switching) oraz zjawisko jitteru wywołane rywalizacją na poziomie sprzętowym.

Dodatkowym aspektem poddanym analizie w tym scenariuszu jest specyfika architektury wdrożenia. Bazy serwerowe (np. Redis, Memcached, Couchbase, Tarantool, FoundationDB) obsługują zapytania poprzez gniazda sieciowe, co pozwala na naturalne skalowanie z wieloma procesami klienckimi, ale obarcza każdą operację narzutem sieciowym. Z kolei bazy osadzone (LevelDB, RocksDB, BerkeleyDB) są ładowane bezpośrednio do przestrzeni adresowej procesora klienta. W ich przypadku badanie wieloprocesowości ujawnia zasadnicze ograniczenia architektoniczne:
\begin{itemize}
    \item \textbf{LevelDB i RocksDB:} Stosują mechanizm wyłącznej blokady (LOCK) na poziomie katalogu bazy danych. Uruchomienie drugiego procesu klienckiego próbującego uzyskać dostęp do zapisu w tej samej bazie danych skutkuje natychmiastowym błędem blokady plików, co uniemożliwia współbieżny dostęp wieloprocesowy. Wprawdzie silnik RocksDB oferuje teoretyczną możliwość jednoczesnego otwarcia bazy przez wiele procesów w trybie tylko do odczytu (ang. \emph{read-only}), jednak w kontekście realizowanych badań wariant ten jest niefunkcjonalny. Wynika to z faktu, że cykl testowy wymaga uprzedniego zasilenia bazy danymi (operacje zapisu), a scenariusze benchmarkowe obejmują również modyfikacje i usuwanie rekordów. W związku z tym pomiary dla konfiguracji wieloprocesowych nie mogą zostać przeprowadzone dla tych silników.
    \item \textbf{BerkeleyDB:} Jako jedyna z badanych bibliotek osadzonych udostępnia mechanizmy synchronizacji i współdzielenia struktur danych w ramach środowiska bazodanowego, co pozwala na współbieżny dostęp wielu procesom klienckim. Eksperymenty mają na celu zmierzenie narzutu na synchronizację międzyprocesową w porównaniu do komunikacji sieciowej baz serwerowych.
\end{itemize}

\section{Mierzone metryki}
Podczas każdego eksperymentu rejestrowane są następujące metryki:
\begin{itemize}
    \item \textbf{Przepustowość (Throughput):} liczba poprawnie obsłużonych operacji na sekundę (ops/sec).
    \item \textbf{Opóźnienie (Latency):} czas wykonania pojedynczego zapytania wyrażony w milisekundach, z uwzględnieniem wartości średniej oraz centyli (p50, p95, p99).
    \item \textbf{Zużycie zasobów (Resource Utilization):} procentowe zużycie procesora (CPU), zajętość pamięci RAM oraz natężenie operacji wejścia/wyjścia na dysk (Block I/O) dla kontenera bazy danych (serwer) oraz kontenera aplikacji benchmarkującej (klient).
\end{itemize}

Dane o zużyciu zasobów (CPU, RAM, Block I/O) rejestrowane są za pomocą polecenia \texttt{docker stats}, które domyślnie próbkuje stan kontenerów z częstotliwością jednej próbki na sekundę. W warunkach intensywnego obciążenia systemu operacyjnego lub dużego ruchu wejścia/wyjścia mechanizm ten może pominąć pojedyncze próbki, co skutkuje lokalnymi przerwami w szeregach czasowych. Zjawisko to traktowane jest jako ograniczenie infrastrukturalne narzędzia pomiarowego i nie wpływa na wiarygodność pozostałych mierzonych metryk. Ze względu na to, że wskaźnik obciążenia dysku (Block I/O) udostępniany przez demona Docker ma charakter skumulowany, dane te są przetwarzane i prezentowane dwojako: jako różnica (delta) między wartością maksymalną a minimalną z próbkowania dla danego testu (pokazująca całkowity wolumen operacji wejścia/wyjścia), oraz na wykresach przebiegu czasowego jako skumulowany przyrost od momentu rozpoczęcia testu. W celu zapewnienia czytelności i odpowiedniej wartości analitycznej wizualizacji, wykresy przebiegu zużycia zasobów w czasie są generowane wyłącznie dla konfiguracji testowych, w których zebrano co najmniej 30 próbek pomiarowych. W przypadku krótszych testów, dla których zebrano mniej niż 30 próbek z narzędzia \texttt{docker stats}, proces generowania wykresu jest automatycznie pomijany ze względu na zbyt małą liczbę punktów pomiarowych do sporządzenia reprezentatywnego profilu obciążenia w czasie. Zestawienie konfiguracji testowych, które nie spełniły tego wymagania i zostały pominięte przy generowaniu wykresów zużycia zasobów w czasie, przedstawiono w tabeli~\ref{tab:resource_samples_skipped}.

% Tabela wygenerowana automatycznie przez skrypt plots/process_results.py
\begin{table}[ht]
\centering
\resizebox{\textwidth}{!}{%
\begin{tabular}{llcccc}
\toprule
Baza danych & System operacyjny & Rozmiar ładunku & Typ obciążenia & Limit CPU & Próbki \\
\midrule
LEVELDB & UBUNTU & 128 B & 1M & 4 CPU & 5 \\
LEVELDB & UBUNTU & 4 KB & 100k & 4 CPU & 3 \\
LEVELDB & WSL & 128 B & 1M & 4 CPU & 4 \\
LEVELDB & WSL & 4 KB & 100k & 4 CPU & 7 \\
MEMCACHED & UBUNTU & 128 B & 1M & 4 CPU & 12 \\
MEMCACHED & UBUNTU & 4 KB & 100k & 4 CPU & 4 \\
MEMCACHED & WSL & 128 B & 1M & 4 CPU & 11 \\
MEMCACHED & WSL & 4 KB & 100k & 4 CPU & 4 \\
REDIS & UBUNTU & 4 KB & 100k & 4 CPU & 19 \\
ROCKSDB & UBUNTU & 128 B & 1M & 4 CPU & 5 \\
ROCKSDB & UBUNTU & 4 KB & 100k & 4 CPU & 2 \\
ROCKSDB & WSL & 128 B & 1M & 4 CPU & 4 \\
ROCKSDB & WSL & 4 KB & 100k & 4 CPU & 4 \\
TARANTOOL & UBUNTU & 4 KB & 100k & 4 CPU & 18 \\
\bottomrule
\end{tabular}%
}
\caption{Konfiguracje testowe niespełniające kryterium minimalnej liczby 30 próbek dla wykresów zużycia zasobów}
\label{tab:resource_samples_skipped}
\end{table}


\section{Implementacja i optymalizacja procedur benchmarkowych}

W celu wiernego odzwierciedlenia cech wydajnościowych poszczególnych systemów bazodanowych, kod aplikacji testowej oraz konfiguracja środowiska kontenerowego zostały zoptymalizowane pod kątem eliminacji wąskich gardeł niezwiązanych z samymi silnikami bazodanowymi.

\subsection{Optymalizacja konfiguracji serwera Memcached}
W domyślnej konfiguracji serwera Memcached wielkość pojedynczego rekordu (klucz wraz z wartością) jest ograniczona do 1~MB, co mogłoby uniemożliwić realizację scenariuszy zaawansowanych przy średnim rozmiarze rekordu. Ponadto domyślne mechanizmy zwalniania pamięci (LRU) mogłyby usuwać wcześniej zapisane dane pod wpływem ciśnienia pamięciowego, zaburzając wyniki testów odczytu i modyfikacji. W celu eliminacji tych ograniczeń, kontener serwera Memcached uruchomiono z następującymi parametrami:
\begin{itemize}
    \item Zwiększenie maksymalnego rozmiaru rekordu do 2~MB poprzez flagę \texttt{-I 2m}.
    \item Przydzielenie 4~GB pamięci operacyjnej za pomocą flagi \texttt{-m 4096}.
    \item Wyłączenie mechanizmów automatycznego usuwania kluczy oraz reorganizacji pamięci podręcznej poprzez parametry \texttt{-o no\_lru\_maintainer,no\_lru\_crawler,slab\_automove=0}.
\end{itemize}

\subsection{Obsługa zakleszczeń i konfiguracja transakcji w BerkeleyDB}
Wbudowana baza BerkeleyDB w scenariuszach wieloprocesowych realizuje blokowanie współbieżnego dostępu do stron i rekordów. W warunkach intensywnego obciążenia generowanego przez wiele niezależnych procesów klienckich systematycznie dochodzi do zjawiska rywalizacji o blokady, co skutkuje błędami zakleszczenia (ang. \emph{deadlocks}) i zgłaszaniem wyjątku \texttt{DBLockDeadlockError}.

Aby zapewnić stabilność przebiegu pomiarów bez przedwczesnego przerywania wątków klienckich, w adapterze BerkeleyDB (\texttt{berkeleydb\_test.py}) zaimplementowano dedykowaną funkcję pomocniczą \texttt{\_retry}. Funkcja ta przechwytuje wyjątki zakleszczenia i automatycznie ponawia próbę wykonania danej operacji (odczytu, zapisu lub usuwania) w pętli. Pozwoliło to na ciągły pomiar wydajności bazy przy jednoczesnym uwzględnieniu narzutu na obsługę konfliktów blokowania. Ponadto w konfiguracji środowiska bazy danych (\texttt{DBEnv}) zastosowano flagę \texttt{DB\_TXN\_WRITE\_NOSYNC}, która instruuje silnik, aby zapisywał logi transakcyjne do buforów systemu operacyjnego zamiast synchronicznego wymuszania zapisu na dysk przy każdym zatwierdzeniu (\emph{fsync}).

W celu umożliwienia bezpiecznego, równoległego dostępu wielu procesów klienckich do tej samej fizycznej bazy danych bez ryzyka uszkodzenia struktury plików (ang. \emph{data corruption}), konieczne było skonfigurowanie współdzielonego środowiska BerkeleyDB (\texttt{DBEnv}). Domyślne uruchomienie bazy bez włączenia podsystemów kontroli współbieżności skutkowałoby brakiem koordynacji zapisów, a w konsekwencji nieodwracalnym uszkodzeniem danych. Z tego względu środowisko bazy zostało zainicjalizowane z zestawem flag aktywujących mechanizmy transakcyjne, blokowania, logowania Write-Ahead Log oraz współdzielonej pamięci podręcznej (\texttt{DB\_INIT\_TXN}, \texttt{DB\_INIT\_LOCK}, \texttt{DB\_INIT\_LOG}, \texttt{DB\_INIT\_MPOOL}). Dodatkowo zastosowanie flag \texttt{DB\_REGISTER} oraz \texttt{DB\_RECOVER} zapewnia automatyczną detekcję stanu środowiska i przeprowadzenie procedury naprawczej w przypadku awarii jednego z procesów, zapobiegając pozostawieniu osieroconych blokad i niespójności na dysku.

\subsection{Konfiguracja przestrzeni i procedury składowane w Tarantool}
Serwer Tarantool w swojej konfiguracji (\texttt{init.lua}) został dostosowany do wydajnego przetwarzania in-memory. Przydzielono 4~GB pamięci RAM dla silnika \texttt{memtx} odpowiadającego za struktury danych w pamięci. Ponadto, w celu uniknięcia kosztów transmisji sieciowej i uproszczenia emulacji klienckiej przy złożonych testach (JSON oraz kolejka FIFO), w bazie danych zdefiniowano dedykowane przestrzenie (odpowiednik tabel):
\begin{itemize}
    \item \texttt{kv}: przestrzeń przeznaczona do prostego przechowywania par klucz-wartość (format: klucz tekstowy, wartość tekstowa) z indeksem głównym typu \texttt{TREE} na kluczu.
    \item \texttt{doc}: przestrzeń dedykowana do operacji na dokumentach JSON (format: klucz tekstowy, dokument typu \texttt{map}) z indeksem głównym.
    \item \texttt{queue}: przestrzeń symulująca kolejkę komunikatów FIFO (format: klucz kolejki, unikalny identyfikator numeryczny \texttt{unsigned}, dokument typu \texttt{map}) z indeksem złożonym na kluczu kolejki oraz identyfikatorze wiadomości.
\end{itemize}
Dodatkowo, bezpośrednio w silniku Tarantool zaimplementowano atomowe procedury składowane (Lua), które są wywoływane przez klienta zdalnie:
\begin{itemize}
    \item \texttt{queue\_pop}: wyszukuje minimalny element (FIFO) w przestrzeni kolejki, usuwa go w sposób atomowy i zwraca do klienta.
    \item \texttt{doc\_update\_field}: dokonuje bezpośredniej i atomowej aktualizacji wskazanego pola w dokumencie bez konieczności pobierania i ponownej serializacji całego rekordu po stronie klienta.
    \item \texttt{doc\_increment\_field}: inkrementuje wartość wskazanego pola liczbowego wewnątrz mapy dokumentu w sposób atomowy.
    \item \texttt{doc\_get\_field}: pobiera wartość konkretnego pola z dokumentu JSON bezpośrednio z poziomu serwera bazodanowego.
\end{itemize}
Dzięki zastosowaniu kompilatora LuaJIT w silniku Tarantool, procedury te charakteryzują się minimalnym czasem wykonania i eliminują narzut sieciowy.

\section{Metodyka syntetycznej oceny punktowej}

W celu syntetycznego porównania badanych systemów wprowadzono punktowy system oceny. Pozwala on na ilościowe zestawienie cech architektonicznych, licencyjnych oraz funkcjonalnych poszczególnych rozwiązań, a także bezpośrednie włączenie wyników wydajnościowych do końcowej oceny. Wszystkie kategorie statyczne (architektoniczno-funkcjonalne) są oceniane przy użyciu wag dopasowanych do ich znaczenia projektowego (suma maks. \textbf{1500 punktów}).

\subsection{Kryteria przyznawania punktów za cechy statyczne}

W ramach oceny statycznej wyróżniono pięć głównych obszarów:
\begin{itemize}
    \item \textbf{Model licencjonowania (0 / 50 / 100 pkt):} Określa stopień elastyczności wykorzystania systemu w projektach komercyjnych i własnościowych bez konieczności udostępniania kodu źródłowego aplikacji. Ocena 100 punktów oznacza licencję wysoce permisywną bez ograniczeń typu copyleft (np. BSD, Apache 2.0). Ocena 50 punktów oznacza występowanie ograniczeń typu copyleft (np. AGPLv3) lub restrykcji komercyjnych (np. BSL 1.1). Ocena 0 punktów oznacza licencję komercyjną bądź zamkniętą bez darmowego wariantu produkcyjnego.
    \item \textbf{Skalowanie i replikacja (0 / 250 / 500 pkt):} Ocenia natywne wsparcie systemu dla skalowania poziomego (sharding) oraz wysokiej dostępności (HA). Ocena 500 punktów oznacza wbudowane mechanizmy shardingu oraz automatycznej replikacji z obsługą przełączania awaryjnego (failover). Ocena 250 punktów oznacza ograniczone skalowanie poziome (np. po stronie aplikacji klienckiej lub bez automatycznego shardingu) lub replikację High Availability. Ocena 0 punktów oznacza architekturę czysto osadzoną (embedded), ograniczoną do jednego procesu i systemu plików.
    \item \textbf{Wsparcie dla transakcji ACID (0 / 150 / 300 pkt):} Ocenia poziom gwarancji spójności danych przy jednoczesnych modyfikacjach. Ocena 300 punktów oznacza pełne wsparcie dla transakcji wielokluczowych (w tym rozproszonych) z zachowaniem pełnej izolacji. Ocena 150 punktów oznacza uproszczoną transakcyjność (np. atomowość na poziomie pojedynczego klucza lub grupowanie operacji bez pełnego wycofywania zmian). Ocena 0 punktów oznacza brak mechanizmów transakcyjnych.
    \item \textbf{Dodatkowe funkcjonalności i unikalność (0 / 250 / 500 pkt):} Ocenia dostępność zaawansowanych mechanizmów wykraczających poza proste operacje Key-Value. Ocena 500 punktów oznacza wielomodelowość i bogate API (np. strumienie, wyszukiwanie pełnotekstowe, wbudowane silniki skryptowe). Ocena 250 punktów oznacza średni stopień rozszerzalności (np. Column Families, różne fizyczne metody dostępu). Ocena 0 punktów oznacza czysty silnik Key-Value bez dodatkowych struktur.
    \item \textbf{Trwałość danych (Persistence) (0 / 50 / 100 pkt):} Ocenia domyślne gwarancje zachowania spójności danych przy nagłej awarii zasilania lub restartu systemu. Ocena 100 punktów oznacza domyślny, synchroniczny zapis transakcji lub logu operacji (WAL) na dysk fizyczny hosta przed potwierdzeniem zapisu klientowi. Ocena 50 punktów oznacza konfigurowalną lub asynchroniczną trwałość (zapis w tle lub cykliczne migawki). Ocena 0 punktów oznacza brak mechanizmów trwałości (silnik działający wyłącznie w pamięci RAM).
\end{itemize}

\subsection{Kryteria przyznawania punktów za wydajność}

Z uwagi na to, że każdy test wydajnościowy (np. Insert, Read) jest uruchamiany w różnych konfiguracjach sprzętowo-programowych (kombinacja 1, 2, 4 rdzeni CPU oraz 1, 2, 4, 8 procesów klienckich, co daje 12 przypadków na jeden test), zdecydowano się na zastosowanie skali o mniejszej rozpiętości w celu uniknięcia nadmiernej dominacji punktowej pojedynczych testów.

W każdej z 12 konfiguracji danego testu bazy są szeregowane, a punkty przyznawane są według poniższej skali:
\begin{itemize}
    \item 1. miejsce: 14 pkt
    \item 2. miejsce: 12 pkt
    \item 3. miejsce: 10 pkt
    \item 4. miejsce: 8 pkt
    \item 5. miejsce: 6 pkt
    \item 6. miejsce: 4 pkt
    \item 7. miejsce: 2 pkt
    \item 8. miejsce / awaria / brak testu: 0 pkt
\end{itemize}

Taki podział pozwala precyzyjnie ocenić i nagrodzić bazy, które charakteryzują się stabilną wydajnością oraz dobrą skalowalnością przy wzroście liczby rdzeni CPU oraz współbieżnych procesów klienckich.

Warto podkreślić, że do ostatecznej oceny punktowej (w rankingu końcowym) brane są pod uwagę wyłącznie wyniki wydajnościowe uzyskane w środowisku natywnego systemu Linux (Ubuntu). Wynika to z faktu, że w rzeczywistych wdrożeniach produkcyjnych systemy bazodanowe powinny być uruchamiane bezpośrednio na dedykowanych systemach operacyjnych Linux w celu zapewnienia maksymalnej wydajności i stabilności, podczas gdy środowisko wirtualizacji WSL 2 (Windows Subsystem for Linux 2) ma charakter deweloperski i służy przede wszystkim do lokalnych prac programistycznych.

Dodatkowo, dla zaawansowanych scenariuszy testowych (kolejka FIFO oraz 6 operacji na dokumentach JSON, co daje łącznie 7 scenariuszy) wprowadza się punkty za realizację funkcjonalności:
\begin{itemize}
    \item \textbf{100 pkt:} pomyślne ukończenie testu z użyciem natywnego wsparcia dla danej struktury danych po stronie serwera bazy danych (np. Redis, Tarantool, Couchbase).
    \item \textbf{50 pkt:} pomyślne ukończenie testu z wykorzystaniem programowej emulacji struktur po stronie aplikacji klienckiej w języku Python (np. LevelDB, RocksDB, BerkeleyDB, Memcached, FoundationDB).
    \item \textbf{0 pkt:} brak możliwości przeprowadzenia testu w danej konfiguracji (np. z powodu blokad wieloprocesowych lub całkowitego braku wsparcia).
\end{itemize}

\subsection{Zestawienie punktowe cech statycznych}

Tabela \ref{tab:punktacja_statyczna} przedstawia bazową punktację architektoniczno-funkcjonalną przed uwzględnieniem wyników wydajnościowych (suma maks. 1500 punktów).

\begin{table}[H]
\centering
\resizebox{\textwidth}{!}{%
\begin{tabular}{|l|c|c|c|c|c|c|}
\hline
\textbf{Baza danych} & \textbf{Licencja} & \textbf{Skalowanie} & \textbf{ACID} & \textbf{Unikalność} & \textbf{Trwałość} & \textbf{Suma statyczna} \\ \hline
FoundationDB & 100 & 500 & 300 & 500 & 100 & 1500 \\ \hline
Tarantool & 100 & 500 & 300 & 500 & 50 & 1450 \\ \hline
Redis & 50 & 500 & 300 & 500 & 50 & 1400 \\ \hline
Couchbase & 50 & 500 & 150 & 500 & 50 & 1250 \\ \hline
BerkeleyDB & 50 & 250 & 300 & 250 & 100 & 950 \\ \hline
RocksDB & 100 & 0 & 150 & 250 & 100 & 600 \\ \hline
Memcached & 100 & 250 & 0 & 0 & 0 & 350 \\ \hline
LevelDB & 100 & 0 & 150 & 0 & 100 & 350 \\ \hline
\end{tabular}%
}
\caption{Zestawienie punktowe cech statycznych badanych systemów}
\label{tab:punktacja_statyczna}
\end{table}

\subsection{Uzasadnienie ocen cech statycznych}

Poniżej przedstawiono szczegółowe uzasadnienie przyznanych ocen w kategoriach statycznych dla poszczególnych systemów bazodanowych:

\paragraph{Redis (1400 pkt)}
\begin{itemize}
    \item \textbf{Model licencjonowania (50 pkt):} Przejście na model trójlicencyjny (AGPLv3, RSALv2, SSPLv1) w wersji 8 ograniczyło elastyczność komercyjną ze względu na wymogi copyleft licencji AGPLv3.
    \item \textbf{Skalowanie i replikacja (500 pkt):} Zapewnione jest pełne wsparcie dla klastrowania poziomego (Redis Cluster) z automatycznym shardingiem i replikacją.
    \item \textbf{Wsparcie dla transakcji ACID (300 pkt):} Mimo braku klasycznego mechanizmu wycofywania zmian (rollback), jednowątkowy model przetwarzania poleceń w połączeniu ze skryptami Lua zapewnia pełną atomowość i najwyższy stopień izolacji (Serializable) wykonywania transakcji w pamięci RAM.
    \item \textbf{Dodatkowe funkcjonalności i unikalność (500 pkt):} Oferowany jest wyjątkowo bogaty zestaw wbudowanych struktur (m.in. listy, zbiory, strumienie, moduły RedisJSON i RediSearch).
    \item \textbf{Trwałość danych (50 pkt):} Zastosowano asynchroniczny zapis na dysk w tle za pomocą cyklicznych migawek RDB oraz logu operacji AOF.
\end{itemize}

\paragraph{Memcached (350 pkt)}
\begin{itemize}
    \item \textbf{Model licencjonowania (100 pkt):} Zastosowanie trzyklauzulowej licencji BSD umożliwia dowolną integrację kodu.
    \item \textbf{Skalowanie i replikacja (250 pkt):} Brak wbudowanego mechanizmu replikacji i shardingu po stronie serwera. Skalowanie poziome realizowane jest po stronie klienta za pomocą algorytmów skrótu (consistent hashing).
    \item \textbf{Wsparcie dla transakcji ACID (0 pkt):} Wykazano całkowity brak wsparcia dla operacji transakcyjnych.
    \item \textbf{Dodatkowe funkcjonalności i unikalność (0 pkt):} Klasyczny, prosty cache typu klucz-wartość bez dodatkowych struktur danych.
    \item \textbf{Trwałość danych (0 pkt):} Przechowywanie danych odbywa się wyłącznie w pamięci RAM, bez mechanizmów trwałości dyskowej.
\end{itemize}

\paragraph{LevelDB (350 pkt)}
\begin{itemize}
    \item \textbf{Model licencjonowania (100 pkt):} Biblioteka udostępniana jest na liberalnej licencji BSD.
    \item \textbf{Skalowanie i replikacja (0 pkt):} Silnik wbudowany (embedded), brak obsługi protokołów sieciowych oraz replikacji.
    \item \textbf{Wsparcie dla transakcji ACID (150 pkt):} Zapewniona atomowość zapisu grupy operacji przy użyciu mechanizmu WriteBatch i dziennika WAL.
    \item \textbf{Dodatkowe funkcjonalności i unikalność (0 pkt):} Podstawowy silnik LSM-Tree oferujący jedynie proste operacje zapisu, odczytu i usuwania ciągu bajtów.
    \item \textbf{Trwałość danych (100 pkt):} Domyślny, silny zapis na dysk fizyczny do logu WAL przed potwierdzeniem zapisu klientowi.
\end{itemize}

\paragraph{RocksDB (600 pkt)}
\begin{itemize}
    \item \textbf{Model licencjonowania (100 pkt):} Dualna licencja (Apache 2.0 / GPLv2) pozwala na bezpieczną integrację w projektach komercyjnych.
    \item \textbf{Skalowanie i replikacja (0 pkt):} Silnik wbudowany (embedded) bez wbudowanej obsługi sieci i replikacji.
    \item \textbf{Wsparcie dla transakcji ACID (150 pkt):} Atomowość zapisu realizowana przez mechanizm WriteBatch i dziennik WAL.
    \item \textbf{Dodatkowe funkcjonalności i unikalność (250 pkt):} Rozbudowane mechanizmy optymalizacji (Column Families, Merge Operators, filtry Blooma).
    \item \textbf{Trwałość danych (100 pkt):} Domyślny, silny zapis sekwencyjny do logu WAL na dysk.
\end{itemize}

\paragraph{Couchbase (1250 pkt)}
\begin{itemize}
    \item \textbf{Model licencjonowania (50 pkt):} Wersja społecznościowa podlega licencji BSL 1.1, która ogranicza komercyjne zastosowania typu SaaS.
    \item \textbf{Skalowanie i replikacja (500 pkt):} Zaawansowany automatyczny sharding (vBuckety) oraz replikacja międzycentrowa (XDCR).
    \item \textbf{Wsparcie dla transakcji ACID (150 pkt):} Zapewniono atomowość na poziomie pojedynczego dokumentu JSON.
    \item \textbf{Dodatkowe funkcjonalności i unikalność (500 pkt):} Silnik zapytań SQL++ (N1QL), wbudowane usługi wyszukiwania pełnotekstowego (FTS), analityki oraz zdarzeń (Eventing).
    \item \textbf{Trwałość danych (50 pkt):} Model memory-first, zapytania potwierdzane są po zapisaniu w pamięci RAM, po czym asynchronicznie zapisywane na dysk.
\end{itemize}

\paragraph{Tarantool (1450 pkt)}
\begin{itemize}
    \item \textbf{Model licencjonowania (100 pkt):} Udostępniany na dwuklauzulowej licencji BSD.
    \item \textbf{Skalowanie i replikacja (500 pkt):} Wbudowana replikacja (w tym synchroniczna oparta na protokole Raft) oraz sharding poziomy realizowany przez bibliotekę vshard.
    \item \textbf{Wsparcie dla transakcji ACID (300 pkt):} Pełne transakcje ACID dla silnika Memtx dzięki jednowątkowej architekturze opartej na kooperatywnych włóknach (fibers).
    \item \textbf{Dodatkowe funkcjonalności i unikalność (500 pkt):} Pełnoprawny serwer aplikacji umożliwiający uruchamianie kodu Lua bezpośrednio w procesie bazy danych.
    \item \textbf{Trwałość danych (50 pkt):} Asynchroniczny lub okresowy zapis dziennika transakcji WAL i zrzut migawek na dysk.
\end{itemize}

\paragraph{FoundationDB (1500 pkt)}
\begin{itemize}
    \item \textbf{Model licencjonowania (100 pkt):} Otwarty kod udostępniany na licencji Apache 2.0.
    \item \textbf{Skalowanie i replikacja (500 pkt):} W pełni automatyczny, przezroczysty sharding poziomy z wbudowaną replikacją tolerującą awarie węzłów.
    \item \textbf{Wsparcie dla transakcji ACID (300 pkt):} Globalne transakcje rozproszone o najwyższym poziomie izolacji (Serializable).
    \item \textbf{Dodatkowe funkcjonalności i unikalność (500 pkt):} Architektura oparta na warstwach (Layers), umożliwiająca emulację dokumentów czy relacji grafowych na bazie podstawowego silnika klucz-wartość.
    \item \textbf{Trwałość danych (100 pkt):} Synchroniczny zapis transakcji na dyski wielu replik dziennika przed zatwierdzeniem operacji zapisu.
\end{itemize}

\paragraph{BerkeleyDB (950 pkt)}
\begin{itemize}
    \item \textbf{Model licencjonowania (50 pkt):} Licencja AGPLv3 wymusza otwarcie kodu aplikacji klienckiej lub zakup licencji komercyjnej od firmy Oracle.
    \item \textbf{Skalowanie i replikacja (250 pkt):} Replikacja w trybie High Availability pozwalająca na skalowanie odczytów, brak jednak automatycznego shardingu poziomego.
    \item \textbf{Wsparcie dla transakcji ACID (300 pkt):} Pełne lokalne transakcje ACID z menedżerem blokad na poziomie stron i rekordów.
    \item \textbf{Dodatkowe funkcjonalności i unikalność (250 pkt):} Obsługa czterech różnych struktur fizycznych przechowywania danych (B-Tree, Hash, Queue, Recno).
    \item \textbf{Trwałość danych (100 pkt):} Transakcyjne wymuszenie zapisu logu WAL na dysk hosta.
\end{itemize}
